\documentclass[../main.tex]{subfiles}
\begin{document}

\chapter{キャッシュコヒーレンス}
\lessoninfo{
  video={Lesson 19，動画 1--33},
  textbook={H\&P \cite{hennessy2017} §5.2--5.4},
  keywords={キャッシュコヒーレンス，バススヌーピング，書き込み時更新方式，書き込み時無効化方式，MSI プロトコル，キャッシュ間転送，MOSI プロトコル，MESI プロトコル，MOESI プロトコル，ディレクトリ方式，コヒーレンスミス，真共有，偽共有},
}

第18章の共有メモリ型マルチプロセッサでは，各コアが自分自身のキャッシュを
持ち，すべてのコアが1つの主記憶を共有する．キャッシュが保持しているのは
データの私的なコピーであり，あるコアが自分のコピーに書き込むと，他の
キャッシュのコピーは---そしてライトバック方式のもとでは主記憶も---古い値
を保持したままである．この章では，そのような古いコピーがあるとマルチ
プロセッサが誤った結果を返すことを示し，それを排除する性質であるキャッシュ
コヒーレンスを定義し，それを提供する機構---共有バス上の書き込み時更新方式
と書き込み時無効化方式，MSI プロトコルとその改良である MOSI，MESI，
MOESI，そしてより大きなマシンのためのディレクトリ方式---を順に組み立てる．
最後に，コヒーレンスが引き起こす追加のキャッシュミスを扱う．
\caution{コヒーレンスは一度に1つのメモリ位置についての性質である．
\emph{異なる}位置へのアクセスがどの順序で見えるかはメモリ一貫性であり，
第21章で扱う．}

\section{コヒーレンスの問題}
\label{sec:l19-problem}

各コアが第12章のライトバック方式に従う専用キャッシュを持つ，集中共有
メモリのマルチプロセッサ（\cref{sec:l18-centralized}）を考える．
\source{Lesson 19，動画 1--3；H\&P §5.2．}単一プロセッサではこの方式は
無害である．主記憶がしばらく古い値を保持することはあるが，それを読みうる
唯一のプロセッサは，自分のキャッシュにある新しいコピーを先に見つける
からである．複数のキャッシュがあると，もはやそうはいかない．
\cref{fig:l19-problem} は，初めに 5 を保持するメモリ位置 X を追ったもの
である．\margin{X は1つのワードだが，キャッシュが扱うのはそれを含む
ブロック全体である．x86-64 と大半の Arm コアでは 64 バイト，Apple の M
シリーズでは 128 バイトである（第12章）．}
\begin{enumerate}
  \item コア 0 が X を読む．読み出しはミスし，コア 0 のキャッシュは X を
    含むブロックを主記憶から読み込む．
  \item コア 1 が X を読み，同様にブロックのコピーを読み込む．
  \item コア 0 が X に 8 を書く．書き込みはヒットし，ライトバック方式の
    もとでは，コア 0 のキャッシュ内のコピーだけが変更され，そのコピーは
    ダーティになる．
  \item コア 1 が X を読む．読み出しはコア 1 のキャッシュでヒットし，
    8 が書かれているにもかかわらず 5 を返す．
  \item コア 2 が X を読む．読み出しはミスし，やはり 5 を保持している
    主記憶から供給される．
\end{enumerate}

\begin{figure}[htbp]
  \centering
  % The cache coherence problem: three cores with private write-back caches
% share one main memory.  Core 0 has overwritten X in its cache only; core 1
% reads a stale copy from its cache, core 2 a stale value from memory.
\begin{tikzpicture}[x=1cm,y=1cm, font=\footnotesize\sffamily, >={Stealth[length=1.8mm]},
  core/.style={draw, fill=ln-accent!15, minimum width=2.5cm, minimum height=0.6cm, inner sep=2pt},
  cache/.style={draw, fill=ln-box, minimum width=2.5cm, minimum height=0.6cm, inner sep=2pt},
  big/.style={draw, fill=ln-box, minimum height=0.7cm, inner sep=3pt},
  op/.style={font=\scriptsize, align=center, anchor=south},
  bad/.style={text=ln-caution}]
  \foreach \i/\c in {0/{X $=$ 8\enspace\iflangja{（ダーティ）}{(dirty)}},
                     1/{X $=$ 5},
                     2/{\iflangja{X のコピーなし}{no copy of X}}} {
    \node[core] (c\i) at (\i*3.4,0) {\iflangja{コア}{core}~\i};
    \node[cache] (k\i) at (\i*3.4,-0.85) {\c};
    \draw (c\i.south) -- (k\i.north);
    \draw (k\i.south) -- (\i*3.4,-1.75);
  }
  % operations, numbered in the order they happen
  \node[op] at (c0.north) {\iflangja{(1) X を読む\\(3) X に 8 を書く}%
                                   {(1) read X\\(3) write 8 to X}};
  \node[op] at (c1.north) {\iflangja{(2) X を読む}{(2) read X}\\
    \textcolor{ln-caution}{\iflangja{(4) X を読む: ヒット，5}%
                                   {(4) read X: hit, 5}}};
  \node[op, bad] at (c2.north) {\iflangja{(5) X を読む:\\ミス，主記憶から 5}%
                                        {(5) read X:\\miss, 5 from memory}};
  % bus and memory
  \draw[line width=1.4pt] (-1.45,-1.75) -- (8.25,-1.75)
    node[right, font=\footnotesize] {\iflangja{バス}{bus}};
  \node[big, minimum width=4.2cm] (mem) at (3.4,-2.75) {\iflangja{主記憶: X $=$ 5}{main memory: X $=$ 5}};
  \draw (mem.north) -- (3.4,-1.75);
\end{tikzpicture}

  \caption{コヒーレントでないキャッシュ．コア 0 とコア 1 が，5 を保持する
    X を読み（手順 1 と 2），コア 0 が自分のキャッシュに 8 を書く
    （手順 3）．その後のコア 1 の読み出し（手順 4）とコア 2 の読み出し
    （手順 5）は，どちらも古い値 5 を返す．}
  \label{fig:l19-problem}
\end{figure}

2つの誤った読み出しの原因は異なる．手順 5 は主記憶から古い値を読む．
書き込みのたびに主記憶を更新するライトスルー方式であれば，これは避け
られる．手順 4 は他のキャッシュにある古いコピーを読んでおり，個々の
キャッシュの書き込み方式ではこれを避けられない．コア 0 のキャッシュが
書き込みをどう扱おうと，コア 1 のキャッシュ内のコピーは変わらないから
である．\caution{ライトスルー方式は専用キャッシュをコヒーレントにしない．
他のキャッシュ内のコピーはやはり古くなる．}

\section{コヒーレンス}
\label{sec:l19-definition}

共有メモリのプログラムは，すべてのコアが共有する，キャッシュのない1つの
メモリの振る舞いを期待する．\source{Lesson 19，動画 4--5；H\&P §5.2．}

\begin{definition}[キャッシュコヒーレンス]
  コアが専用キャッシュを持つマルチプロセッサは，すべてのメモリ位置 X に
  ついて次が成り立つとき，\term[きゃっしゅこひーれんす]{キャッシュ%
  コヒーレンス}[cache coherence]を提供するという．
  \begin{enumerate}
    \item あるコアでの X の読み出しは，その間に他のコアが X に書いて
      いなければ，そのコアが X に最後に書いた値を返す．
    \item あるコアが X に書き，十分に後に別のコアが X を読み，その間に
      X への書き込みがなければ，その読み出しは最初のコアが書いた値を
      返す．
    \item X への書き込みは直列化される．すなわち，すべてのコアが X への
      任意の2つの書き込みを同じ順序で観測する．
  \end{enumerate}
\end{definition}

第1の条件は単一プロセッサにおけるメモリの振る舞いであり，どのキャッシュ
もすでにこれを保っている．第2の条件はそれを他のコアへ拡張するが，遅延は
許す．別のコアの書き込みがまだ進行中の間に発行された読み出しは，古い値を
返してよい．第3の条件により，ほぼ同時の2つの書き込みのどちらが後だったか
について，すべてのコアの見解が一致する．\cref{fig:l19-problem} の手順 4
と 5 の読み出しは第2の条件に違反している．すなわち，これらのキャッシュは
\emph{コヒーレントでない}．

\begin{note}
  コヒーレンスは，読み出しが何を返すかに対する要請であって，コピーが何を
  保持しているかに対する要請ではない．コピーどうしが食い違っていても
  コヒーレンスに違反するとは限らない．主記憶はキャッシュ内のダーティな
  ブロックより古い値を保持しうるし，キャッシュも，無効と印が付いたライン
  に古い値を残していてよい．どの読み出しも最新の値を保持するコピーから
  供給される限り，システムはコヒーレントである．
\end{note}

\section{コヒーレンスの実現}
\label{sec:l19-achieving}

定義と \cref{fig:l19-problem} の例は，コヒーレンスの機構が何をしなければ
ならないかを示している．\source{Lesson 19，動画 6；H\&P §5.2．}
\begin{enumerate}
  \item \textbf{書き込みを可視にする．}ライトバック方式では，ヒットする
    書き込みはコアの外に何の動作も生じさせないので，他のキャッシュは
    自分のコピーが古くなったことを知りえない．したがって，他のキャッシュ
    が保持しているかもしれないブロックへの書き込みは，それらに知らせ
    なければならない．
  \item \textbf{他のコアの書き込みに対処する．}他のコアが書いたブロックの
    コピーを保持するキャッシュは，自分のコピーを新しい値で\emph{更新}する
    か，\emph{無効化}して次のアクセスがミスし新しい値を取得するように
    しなければならない．
  \item \textbf{書き込みを順序付ける．}異なるコアによる同じブロックへの
    書き込みは，すべてのキャッシュが観測する1つの順序で，すべてのコピーに
    適用されなければならない．
  \item \textbf{読み出しを最新のコピーから供給する．}ミスした読み出しは
    最新の値を得なければならないが，その値は主記憶ではなく他のキャッシュ
    にあるかもしれない．
\end{enumerate}

集中共有メモリでは，すべてのキャッシュと主記憶が1本のバスに接続され，
バス上のあらゆる操作が，接続されたすべての装置に届く．そのため
コヒーレンスは\term[ばすすぬーぴんぐ]{バススヌーピング}[bus snooping]の
上に構築できる．各キャッシュはバス上のすべての操作を---主記憶宛てや他の
キャッシュ宛てのものも含めて---監視し，自分が保持するブロックに関わる
ものに反応する．バスは一度に1つの操作しか運ばず，すべてのキャッシュが
その順序で操作を観測するので，スヌーピングは操作の順序付けも行う．共有
バスを持たないマシンは，この2つの仕事のために別の機構，すなわち
\cref{sec:l19-directory} のディレクトリを必要とする．このように
コヒーレンスの機構は，他のコピーを更新するか無効化するか，そして書き込み
をスヌーピングで観測・順序付けするかディレクトリでそうするか，という2つの
選択の組合せである．

\section{書き込み時更新方式}
\label{sec:l19-write-update}

書き込みに対処する第1の方法は，新しい値をすべてのコピーへ伝播させること
である．\source{Lesson 19，動画 7--11；H\&P §5.2．}

\begin{definition}[書き込み時更新方式]
  \term[かきこみじこうしんほうしき]{書き込み時更新方式}%
  [write-update coherence]では，すべての書き込みがそのアドレスと新しい値
  とともにブロードキャストされ，書かれたブロックのコピーを保持する他の
  すべてのキャッシュが自分のコピーを更新する．
\end{definition}

スヌーピングのもとでは，各コアはすべての書き込みをそのままバスへ出す．
主記憶は他の書き込みと同様にそれを実行し，他のキャッシュはそれをスヌープ
して，当該ブロックを保持していれば自分のコピーを更新する．
\caution{スヌーピングは他のキャッシュを探しに行くことではない．各
キャッシュはバスを監視し，自分のタグだけを引く．}
\cref{fig:l19-problem} では，手順 3 の書き込みが主記憶とコア 1 の
キャッシュの X を 8 に変え，その後の2つの読み出しはどちらも 8 を返す
ことになる．バスは一度に1つの操作しか運ばないので，2つのコアが同じ位置に
同時に書いた場合でも，異なるコアによる書き込みは同じ順序ですべてのコピー
に適用される．これはコヒーレンスの第3の条件が要求するとおりである．

この基本的な方式は正しいが高価である．すべての書き込みが共有バスを占有し，
さらに主記憶でも実行される---これはまさにライトバック方式が避けようとする
コストである．2つの最適化がそのほとんどを取り除く．

\subsection{主記憶への書き込みを避ける}

第12章のダーティビットを使えば，キャッシュは主記憶への書き込みを先送り
できる．書き込みは，他のキャッシュがコピーを更新できるようにバスへ出され
るが，主記憶はそれを無視する．代わりに，書き込んだキャッシュがブロックの
ダーティビットを立て，そのブロックを追い出すときに主記憶へ書き戻す．
自分がダーティに保持するブロックへの他のコアの書き込みをスヌープした
キャッシュは，自分のコピーを更新してダーティビットを下ろす．こうして，
最後に書いたキャッシュだけがそのブロックに責任を持つ．ダーティなコピーが
存在する間，主記憶は古い．そのためダーティなコピーを持つキャッシュは，
そのブロックの\emph{読み出し}についてもバスをスヌープする．他のコアが
そのブロックでミスして読み出しをバスへ出したとき，このキャッシュがデータ
を供給し，主記憶は応答しない．

\subsection{バスへの書き込みを避ける}

大半のブロックは1つのコアだけが使うので，そのようなブロックへの書き込みを
ブロードキャストしてもバスを無駄に占有するだけである．各キャッシュライン
の\emph{共有ビット}は，他のキャッシュがそのブロックを保持している可能性
があるかどうかを記録する．\margin{共有の信号を運ぶのはワイヤードOR の
1本の線である．ブロックを保持するキャッシュがそれを引くので，すべての
キャッシュが同時に答える．}キャッシュは，自分が保持するブロックの
読み出しをスヌープすると，自分の共有ビットを立て，コピーが存在することを
バスの\emph{共有線}上で知らせる．読み出しを行ったキャッシュは，この信号が
あれば自分の共有ビットを立てる．共有ビットが下りているブロックへの
書き込みは，単一プロセッサと同様に局所的に行われ，共有ビットが立っている
ブロックへの書き込みだけがバスへ出される．\margin{共有ビットは保守的で
ある．他のコピーが追い出された後も立ったままになる．}

\begin{example}
  \label{ex:l19-write-update}
  2つのコアが，初めはどのキャッシュも保持していないブロック X にアクセス
  する．\cref{tab:l19-write-update} は，6回のアクセスについて，基本方式の
  場合と最適化を加えた場合のバス操作数と主記憶への書き込み数を数えたもの
  である．ダーティビットがあれば，この系列の間に主記憶は書かれず（コア 0
  が追い出すときに一度だけ書き戻す），手順 4 の読み出しはコア 0 が供給
  する．さらに共有ビットがあれば，コア 1 がコピーを持つ前の2つの書き込み
  は局所的なままとなり，バスへ出されるのは手順 5 の書き込みだけである．
\end{example}

\begin{table}[htbp]
  \centering\small
  \caption{書き込み時更新方式（\cref{ex:l19-write-update}）．rd: バス上の
    読み出し．wr: バス上の書き込み，または主記憶への書き込み．rd$^{*}$:
    主記憶ではなくコア 0 が供給した読み出し．}
  \label{tab:l19-write-update}
  \setlength{\tabcolsep}{4pt}
  \begin{tabular}{@{}cl cc cc cc@{}}
    \toprule
    & & \multicolumn{2}{c}{基本方式} & \multicolumn{2}{c}{ダーティビット}
      & \multicolumn{2}{c}{\shortstack{ダーティビットと\\共有ビット}} \\
    \cmidrule(lr){3-4}\cmidrule(lr){5-6}\cmidrule(l){7-8}
    手順 & アクセス & バス & 主記憶 & バス & 主記憶 & バス & 主記憶 \\
    \midrule
    1 & コア 0 が X を読む  & rd  & --- & rd       & --- & rd       & --- \\
    2 & コア 0 が X に書く & wr  & wr  & wr       & --- & ---      & --- \\
    3 & コア 0 が X に書く & wr  & wr  & wr       & --- & ---      & --- \\
    4 & コア 1 が X を読む  & rd  & --- & rd$^{*}$ & --- & rd$^{*}$ & --- \\
    5 & コア 0 が X に書く & wr  & wr  & wr       & --- & wr       & --- \\
    6 & コア 1 が X を読む  & --- & --- & ---      & --- & ---      & --- \\
    \midrule
    \multicolumn{2}{@{}l}{合計} & 5 & 3 & 5 & 0 & 3 & 0 \\
    \bottomrule
  \end{tabular}
\end{table}

\section{書き込み時無効化方式}
\label{sec:l19-write-invalidate}

書き込みに対処する第2の方法は，他のコピーを更新する代わりに取り除くこと
である．\source{Lesson 19，動画 12--15；H\&P §5.2．}

\begin{definition}[書き込み時無効化方式]
  \term[かきこみじむこうかほうしき]{書き込み時無効化方式}%
  [write-invalidate coherence]では，他のキャッシュが保持しているかも
  しれないブロックに書き込むコアが，まずそのブロックの無効化をブロード
  キャストし，そのブロックのコピーを保持する他のすべてのキャッシュが，
  有効ビットを下ろしてコピーを破棄する．
\end{definition}

無効化はブロックのアドレスだけを運ぶ．いったんブロードキャストされれば，
書き込んだキャッシュが唯一の有効なコピーを保持するので，そのブロックへの
---ブロック内のどのアドレスへのものであっても---その後の書き込みは
局所的であり，別のコアがそのブロックにアクセスするまでバス操作を必要と
しない．そのアクセスは，コピーが無効化されているのでミスし，ブロックが
書き戻されていれば主記憶から，そうでなければそれを書いたキャッシュから
最新の値を得る．

\subsection{更新と無効化の比較}

どちらの方式のバス操作が少なくて済むかは，アクセスのパターンに依存する．
\begin{itemize}
  \item \textbf{別のコアが読む前に1つのブロックへ複数回書き込む場合．}
    書き込み時更新方式はすべての書き込みをブロードキャストするが，
    書き込み時無効化方式は1回の無効化をブロードキャストするだけである．
    無効化はブロック単位で働くので，これは1つのアドレスへの繰り返しの
    書き込みについても，同じブロック内の異なるアドレスへの書き込みに
    ついても同様に成り立つ．
  \item \textbf{スレッドが別のコアへ移る場合．}そのスレッドが使っていた
    ブロックは元のコアのキャッシュに残り，スレッドが新しいコアでそれらを
    読むと，どちらのキャッシュもそれらを共有と記録する．書き込み時更新
    方式では，そのようなブロックへのスレッドのその後の書き込みはすべて
    ブロードキャストされ，元のコアがもう使わないコピーを，そのコピーが
    たまたま追い出されるまで更新し続ける．書き込み時無効化方式では，最初の
    書き込みが古いコピーを取り除き，その後の書き込みはすべて局所的で
    ある．
  \item \textbf{異なるコアによる書き込みと読み出しが交互に起こる場合．}
    あるコアがブロックに書き，別のコアが書き込みのたびに新しい値を読む
    なら，書き込み時更新方式は書き込みごとに1回のバス書き込みを必要と
    し，読み出しはすべてヒットする．書き込み時無効化方式は，書き込み
    ごとに無効化を，読み出しごとにミスを必要とする．
\end{itemize}
書き込み時更新方式が有利なのは最後のパターンだけであり，最初の2つのほうが
一般的な場合である．\margin{書き込み時更新方式がハードウェアになったのは 1980 年代の少数のマシン---Xerox PARC の Dragon
と DEC SRC の Firefly---だけであり，以後のプロセッサはすべて無効化する．}%
したがってマルチプロセッサは書き込み時無効化方式を
用いる．この章の残りのプロトコルはすべてこれに基づく．

\begin{example}
  コア 0 とコア 1 がともにブロック X のコピーを保持しているとする．コア 0
  が X に 10 回書き，その後コア 1 が X を 1回読むなら，書き込み時更新
  方式は 10 回のバス書き込みを必要とし，コア 1 の読み出しはヒットする．
  書き込み時無効化方式は，1回の無効化とコア 1 のミスのための1回のバス
  読み出し，合計 2回のバス操作で済む．逆にコア 1 が 10 回の書き込みの
  それぞれの後に X を読むなら，書き込み時更新方式は依然として 10 回のバス
  操作で済むが，書き込み時無効化方式は 10 回の無効化と 10 回のミス，合計
  20 回を必要とする．
\end{example}

\begin{keypoint}
  書き込み時更新方式は，書き込みごとに1回のバス操作を代償として，すべての
  コピーを最新に保つ．書き込み時無効化方式は，ブロックごとに1回の無効化を
  払い，その後は書き込んだコアが局所的に処理を進められるようにする．
  書き込みは通常1つのコアからまとまって来るので，実用的なコヒーレンス
  プロトコルは無効化を行う．
\end{keypoint}

\section{MSI プロトコル}
\label{sec:l19-msi}

スヌーピングによる書き込み時無効化方式は，キャッシュラインごとの小さな
状態機械で記述される．\source{Lesson 19，動画 16--19；H\&P §5.2．}ライン
にあるブロックの状態は，コア自身のアクセスと，キャッシュがバス上で観測
する他のコアのアクセスとによって変化する．そのような状態機械のうち最も
単純な \term[MSI ぷろとこる]{MSI プロトコル}[MSI protocol]は，ラインの
有効ビットとダーティビットを使って3つの状態を区別する．
\begin{description}
  \item[M（modified）] 有効かつダーティ．ブロックはこのキャッシュで
    書き込まれ，他のすべてのコピーは無効化されている．すなわち，この
    キャッシュが唯一の有効なコピーを保持しており，その内容は主記憶の
    ブロックより新しい．読み出しも書き込みも局所的である．
  \item[S（shared）] 有効かつダーティでない．他のキャッシュも有効な
    コピーを保持しているかもしれない．読み出しは局所的だが，書き込みは
    まず無効化をブロードキャストしなければならない．
  \item[I（invalid）] 有効ビットが下りている．どのアクセスもミスする．
\end{description}

\begin{figure}[htbp]
  \centering
  % MSI state machine of one block in one cache.  Solid edges: access by this
% core; dashed edges: access by another core, observed on the bus.
\begin{tikzpicture}[x=1cm,y=1cm, font=\footnotesize\sffamily, >={Stealth[length=1.8mm]},
  st/.style={circle, draw, minimum size=0.9cm, inner sep=1pt, fill=ln-box},
  dirty/.style={st, fill=ln-accent!15},
  lab/.style={font=\scriptsize, inner sep=2pt},
  other/.style={dashed, ln-accent}]
  \node[st]    (I) at (0,0)     {I};
  \node[st]    (S) at (3.4,-1.8) {S};
  \node[dirty] (M) at (6.8,0)   {M};
  % this core
  \draw[->] (I) to[bend right=22] node[lab, below left]
    {\iflangja{R / バス読み出し}{R / bus read}} (S);
  \draw[->] (I) -- node[lab, above]
    {\iflangja{W / バス読み出し＋無効化}{W / bus read + invalidate}} (M);
  \draw[->] (S) to[bend left=22] node[lab, above left, pos=0.3]
    {\iflangja{W / 無効化}{W / invalidate}} (M);
  % other cores
  \draw[->, other] (S) to[bend right=22] node[lab, above right, pos=0.45] {W$'$} (I);
  \draw[->, other] (M) to[bend left=22] node[lab, below right]
    {\iflangja{R$'$ / 書き戻し}{R$'$ / write back}} (S);
  \draw[->, other] (M) to[bend right=30] node[lab, above]
    {\iflangja{W$'$ / 書き戻し}{W$'$ / write back}} (I);
\end{tikzpicture}

  \caption{1つのキャッシュ内の1つのブロックに対する MSI プロトコル．R，W
    はこのコアによる読み出しと書き込み，R$'$，W$'$（破線）はバス上で
    観測される他のコアによる読み出しと書き込みである．辺のないアクセスは
    状態を変えず，バス操作も必要としない．網掛けの状態は，主記憶より
    新しいブロックを保持している．}
  \label{fig:l19-msi}
\end{figure}

\cref{fig:l19-msi} は遷移を示す．コア自身のアクセス（実線の辺）は次のよう
に扱われる．
\begin{itemize}
  \item I での読み出しはミスする．キャッシュは読み出しをバスへ出し，
    ブロックを読み込んで S に入る．
  \item I での書き込みもミスする．キャッシュは1回のバス操作でブロックを
    読み，他のすべてのコピーを無効化して，M に入る．
  \item S での書き込みは，無効化をブロードキャストして M に入る．
  \item S での読み出し，および M での読み出しと書き込みは，バス操作を
    必要としない．
\end{itemize}
他のコアのアクセス（破線の辺）はバス上で観測される．
\begin{itemize}
  \item 他のコアの書き込みは，その無効化によって知らされ，S から I への
    遷移を起こす．主記憶が同じ値を保持しているので，コピーは書き戻され
    ずに破棄される．
  \item M では，キャッシュがそのブロックの唯一の最新の値を保持している．
    他のコアがそのブロックを読むと，キャッシュはブロックを主記憶へ書き
    戻し，他のコアもコピーを持つことになるので S に入る．他のコアが
    そのブロックに書くと，キャッシュはブロックを主記憶へ書き戻し，
    書き込んだコアはそこからブロックを得て，このキャッシュは I に入る．
\end{itemize}
キャッシュがブロックを追い出すとき，主記憶へ書き戻すのは M の場合だけで
あり，S のブロックは単に破棄される．\tip{MSI の辺はすべて3つの規則の
どれかに従う．このコアの書き込みは M へ，他のコアの書き込みは I へ，
読み出しは S へ至る．}

\subsection{キャッシュ間転送}
\label{sec:l19-c2c}

あるコアが，別のキャッシュが M で保持するブロックでミスすると，それが
読み出しであれ書き込みであれ，主記憶はバス上の読み出しに古い値で応答して
しまう．その読み出しをスヌープする M のキャッシュは，これを2つの方法で
防げる．
\begin{enumerate}
  \item \emph{中断と再試行}: 読み出しを中断させ，ブロックを主記憶へ書き
    戻す．読み出しを行ったコアは読み出しを再度発行し，主記憶が現在の値で
    それに応答する．これには主記憶への書き込み1回と主記憶からの読み出し
    1回のコストがかかる．
  \item \emph{介入}: 主記憶が応答してはならないことを知らせ，自分で
    ブロックをバスへ出す．読み出しミスでは，両方のコピーがダーティでない
    S に落ち着くので，主記憶もバスからブロックを取り込む．
\end{enumerate}
介入，すなわち\term[きゃっしゅかんてんそう]{キャッシュ間転送}%
[cache-to-cache transfer]は，より複雑なバスのハードウェアを必要とするが，
1回の転送でミスに応答し，読み出しミスの場合には，MSI が M から S への
遷移で要求する書き戻しも同時に行う．書き込み時更新方式のダーティビットの
最適化（\cref{sec:l19-write-update}）も同じように読み出しに応答するが，
そちらではキャッシュはダーティのままであり，主記憶はブロックを取り込ま
ない．\margin{キャッシュ間の転送のコストは L3 ヒット程度，Skylake で約 40
サイクルである．DRAM アクセスには 60--100~ns かかる（第14章）．}

\begin{example}
  \label{ex:l19-msi}
  3つのコアが，初めはどのキャッシュも保持していないブロック X にアクセス
  する（\cref{tab:l19-msi}）．手順 1 と 2 の読み出しは，S の2つのコピーを
  残す．手順 3 のコア 1 の書き込みはコア 0 のコピーを無効化し，その2回目の
  書き込み（手順 4）は局所的である．手順 5 でコア 2 の読み出しがバスに
  現れたとき，唯一の最新の値を保持しているのはコア 1 である．コア 1 は
  ブロックを主記憶へ書き戻し，両方のコピーは S に落ち着く．手順 6 の
  コア 0 の書き込みはミスし，その両方を無効化する．7回のアクセスのうち
  5回がバス操作を必要とし，1回が主記憶への書き込みを引き起こす．
\end{example}

\begin{table}[htbp]
  \centering\small
  \caption{MSI プロトコルのもとでのブロック X（\cref{ex:l19-msi}）．
    アクセス後の各キャッシュにおけるブロックの状態．}
  \label{tab:l19-msi}
  \setlength{\tabcolsep}{4pt}
  \begin{tabular}{@{}cl ccc l@{}}
    \toprule
    手順 & アクセス & コア 0 & コア 1 & コア 2 & バスと主記憶 \\
    \midrule
    1 & コア 0 が X を読む  & S & I & I & 読み出し，主記憶が供給 \\
    2 & コア 1 が X を読む  & S & S & I & 読み出し，主記憶が供給 \\
    3 & コア 1 が X に書く & I & M & I & 無効化 \\
    4 & コア 1 が X に書く & I & M & I & --- \\
    5 & コア 2 が X を読む  & I & S & S & 読み出し，コア 1 が X を書き戻す \\
    6 & コア 0 が X に書く & M & I & I & 読み出しと無効化 \\
    7 & コア 0 が X を読む  & M & I & I & --- \\
    \bottomrule
  \end{tabular}
\end{table}

\section{MOSI プロトコル}
\label{sec:l19-mosi}

キャッシュ間転送は，M から S への遷移における主記憶への書き込みを取り除き
はしない．\source{Lesson 19，動画 20--21．}転送の後，両方のコピーは
ダーティでない S にあるので，どちらのキャッシュも追い出しの際にコピーを
破棄してしまう．転送の間に主記憶がブロックを取り込まない限り，両方の
コピーがなくなった時点で最新の値は失われる．したがって，あるコアが書き，
別のコアが読むことを交互に繰り返すブロックは，どちらのキャッシュもそれを
追い出さないにもかかわらず，書き込みごとに主記憶への書き込みのコストを
払うことになる．

\term[MOSI ぷろとこる]{MOSI プロトコル}[MOSI protocol]は，共有された
ブロックのコピーの1つをその所有者とすることで，これらの書き込みを避ける．
\begin{description}
  \item[O（owned）] 有効．他のキャッシュが S でコピーを保持しているかも
    しれない．ブロックは主記憶のブロックより新しく，このキャッシュが
    それに責任を持つ．すなわち，他のコアによるそのブロックの読み出しに
    応答し，追い出す際にブロックを主記憶へ書き戻す．読み出しは局所的で
    あり，書き込みは無効化をブロードキャストして M に入る．
    \caution{\enquote{所有}は\enquote{排他}ではない．他のキャッシュは S で有効なコピーを保持し
    うる．所有者とは，そのブロックに責任を持つキャッシュのことにすぎ
    ない．}
\end{description}

\begin{figure}[htbp]
  \centering
  % MOSI state machine: MSI plus the owned state O.  Solid edges: access by
% this core; dashed edges: access by another core, observed on the bus.
\begin{tikzpicture}[x=1cm,y=1cm, font=\footnotesize\sffamily, >={Stealth[length=1.8mm]},
  st/.style={circle, draw, minimum size=0.9cm, inner sep=1pt, fill=ln-box},
  dirty/.style={st, fill=ln-accent!15},
  lab/.style={font=\scriptsize, inner sep=2pt},
  other/.style={dashed, ln-accent}]
  \node[st]    (I) at (0,0)      {I};
  \node[st]    (S) at (3.4,-1.8) {S};
  \node[dirty] (M) at (6.8,0)    {M};
  \node[dirty] (O) at (9.8,0)    {O};
  % this core
  \draw[->] (I) to[bend right=22] node[lab, below left]
    {\iflangja{R / バス読み出し}{R / bus read}} (S);
  \draw[->] (I) -- node[lab, above]
    {\iflangja{W / バス読み出し＋無効化}{W / bus read + invalidate}} (M);
  \draw[->] (S) to[bend left=22] node[lab, above left, pos=0.3]
    {\iflangja{W / 無効化}{W / invalidate}} (M);
  \draw[->] (O) to[bend left=30] node[lab, below]
    {\iflangja{W / 無効化}{W / invalidate}} (M);
  % other cores
  \draw[->, other] (S) to[bend right=22] node[lab, above right, pos=0.45] {W$'$} (I);
  \draw[->, other] (M) to[bend right=30] node[lab, above]
    {\iflangja{W$'$ / ブロックを供給}{W$'$ / supply block}} (I);
  \draw[->, other] (M) to[bend left=30] node[lab, above]
    {\iflangja{R$'$ / ブロックを供給}{R$'$ / supply block}} (O);
  \draw[->, other] (O) to[out=-60, in=-120, looseness=5] node[lab, below]
    {\iflangja{R$'$ / ブロックを供給}{R$'$ / supply block}} (O);
  \draw[->, other, rounded corners=6pt] (O.north) -- ++(0,1.75) -| node[lab, above, pos=0.25]
    {W$'$} (I.north);
\end{tikzpicture}

  \caption{MOSI プロトコル．記法は \cref{fig:l19-msi} と同じ．網掛けの
    状態は，主記憶より新しいブロックに責任を持つ．O での W$'$ では，
    書き込んだコアがミスした場合にのみブロックが供給される．}
  \label{fig:l19-mosi}
\end{figure}

キャッシュが M で保持するブロックを他のコアが読むと，キャッシュは
キャッシュ間転送でブロックを供給し，主記憶に書き込むことなく，S ではなく
O に入る（\cref{fig:l19-mosi}）．O では，その後のすべての読み出しに対して
ブロックを供給する．所有者による書き込みは無効化をブロードキャストして M
に戻る．他のコアによる書き込みは O から I への遷移を起こすが，所有者が
まずブロックを供給するのは，書き込んだコアがミスした場合だけである．S に
いる書き込み側はすでに現在の値を保持しているからである．ブロックが主記憶
に達するのは，キャッシュが M または O でそれを追い出すときだけである．
\margin{O 状態は AMD のものである．AMD64 のマニュアルは MOESI を規定して
おり，ダーティなラインは主記憶を経ずにコアからコアへ渡れる．}
\cref{ex:l19-msi} では，MOSI なら手順 5 でコア 1 は S ではなく O に入り，
主記憶への書き込みは起こらず，手順 6 のコア 0 の書き込みに対してはコア 1
がブロックを供給することになる．

\section{排他状態}
\label{sec:l19-mesi}

MSI と MOSI は，非常によくあるアクセスパターンを非効率に扱う．
\source{Lesson 19，動画 22--25；H\&P §5.2；
\cite{papamarcos1984}（Illinois プロトコル）．}
プログラムはしばしば，インクリメントのように，変数を読んでから書く．
そして大半の変数は1つのスレッドだけが使う．そのようなブロックでは，
読み出しがミスして S に入り，続く書き込みが---他のどのキャッシュもコピー
を保持していないにもかかわらず---無効化をブロードキャストする．
このような読み出しと書き込みの組は，単一プロセッサならミスだけで済む
ところ，2回のバス操作を必要とする．

S は，キャッシュには区別できない2つの場合を1つにまとめている．他の
キャッシュがコピーを保持している場合と，ブロックが共有されておらず
まだ書かれていない場合である．第4の状態がこれらを分ける．
\begin{description}
  \item[E（exclusive）] 有効かつダーティでない．他のどのキャッシュも
    コピーを保持していない．読み出しは局所的であり，書き込みも局所的で，
    バス操作なしに M へ入る．\tip{2つの問いで状態が決まる．他の
    キャッシュがコピーを持ちうるか，主記憶が古いか．M はいいえ・はい，
    E はいいえ・いいえ，S ははい・いいえ，O ははい・はいである．}
\end{description}
I での読み出しは，他のどのキャッシュもそのブロックを保持していなければ E
に，そうでなければ S に入る．キャッシュは，どちらの場合かをバス読み出し
そのものの間に知る．ブロックを保持するキャッシュがその読み出しをスヌープ
し，\cref{sec:l19-write-update} と同様にバスの共有線上で知らせるから
である．E で保持しているブロックを他のコアが読むと，主記憶が最新なので，
キャッシュは書き込みなしに S に入る．他のコアによる書き込みは E から I
への遷移を起こす．

MSI に E を加えたものが \term[MESI ぷろとこる]{MESI プロトコル}%
[MESI protocol]であり，MOSI に加えたものが
\term[MOESI ぷろとこる]{MOESI プロトコル}[MOESI protocol]である
（\cref{fig:l19-moesi}）．\cref{tab:l19-states} は5つの状態をまとめた
ものである．

\begin{figure}[htbp]
  \centering
  % MOESI state machine: MOSI plus the exclusive state E.  Solid edges: access
% by this core; dashed edges: access by another core, observed on the bus.
% A write by another core takes every valid state to I (edges not drawn).
\begin{tikzpicture}[x=1cm,y=1cm, font=\footnotesize\sffamily, >={Stealth[length=1.8mm]},
  st/.style={circle, draw, minimum size=0.9cm, inner sep=1pt, fill=ln-box},
  dirty/.style={st, fill=ln-accent!15},
  lab/.style={font=\scriptsize, inner sep=2pt},
  other/.style={dashed, ln-accent}]
  \node[st]    (I) at (0,0)      {I};
  \node[st]    (E) at (3.4,1.8)  {E};
  \node[st]    (S) at (3.4,-1.8) {S};
  \node[dirty] (M) at (6.8,0)    {M};
  \node[dirty] (O) at (9.8,0)    {O};
  % this core
  \draw[->] (I) -- node[lab, above, sloped]
    {\iflangja{R，他にコピーなし}{R, no other copy}} (E);
  \draw[->] (I) -- node[lab, below, sloped]
    {\iflangja{R，他にコピーあり}{R, other copies}} (S);
  \draw[->] (I) -- node[lab, above, pos=0.74, align=center]
    {\iflangja{W / バス読み出し\\＋無効化}{W / bus read\\+ invalidate}} (M);
  \draw[->] (E) -- node[lab, above, sloped] {W} (M);
  \draw[->] (S) -- node[lab, below, sloped]
    {\iflangja{W / 無効化}{W / invalidate}} (M);
  \draw[->] (O) to[bend left=30] node[lab, below]
    {\iflangja{W / 無効化}{W / invalidate}} (M);
  % other cores
  \draw[->, other] (E) -- node[lab, right, pos=0.25] {R$'$} (S);
  \draw[->, other] (M) to[bend left=30] node[lab, above]
    {\iflangja{R$'$ / ブロックを供給}{R$'$ / supply block}} (O);
  \draw[->, other] (O) to[out=-60, in=-120, looseness=5] node[lab, below]
    {\iflangja{R$'$ / ブロックを供給}{R$'$ / supply block}} (O);
\end{tikzpicture}

  \caption{MOESI プロトコル．記法は \cref{fig:l19-msi} と同じ．I から出る
    すべての遷移はバス読み出しを行う．他のコアによる書き込みは，すべての
    有効な状態を I へ移す．書き込んだコアがミスした場合は，M または O の
    キャッシュがブロックを供給する（これらの辺は描いていない）．}
  \label{fig:l19-moesi}
\end{figure}

\begin{table}[htbp]
  \centering\small
  \caption{MOESI の状態．MSI は M，S，I を使い，MOSI は O を，MESI は E を
    加える．責任を持つキャッシュは，他のコアによる読み出しに応答し，
    追い出しの際にブロックを書き戻す．}
  \label{tab:l19-states}
  \setlength{\tabcolsep}{4pt}
  \begin{tabular}{@{}lcccc@{}}
    \toprule
    状態 & 有効 & 他のコピー & 責任 & \shortstack{バスなしの\\書き込み} \\
    \midrule
    M（modified）  & はい   & なし     & はい   & はい \\
    O（owned）     & はい   & ありうる & はい   & いいえ \\
    E（exclusive） & はい   & なし     & いいえ & はい \\
    S（shared）    & はい   & ありうる & いいえ & いいえ \\
    I（invalid）   & いいえ & ありうる & いいえ & いいえ \\
    \bottomrule
  \end{tabular}
\end{table}

\needspace{6\baselineskip}
\begin{example}
  \label{ex:l19-protocols}
  2つのコアが，初めはどのキャッシュも保持していないブロック X にアクセス
  する．コア 0 が X を読んでから書き，コア 1 がそれを読み，コア 0 が
  もう一度書き，コア 1 がもう一度読む．\margin{Intel のソケット間
  プロトコルは第5の状態 F（forward）を加える．O と同様に読み出しに応答する
  共有者を1つ定めるが，F はダーティでない．}\cref{tab:l19-protocols} は，4つの
  プロトコルのもとでの状態と操作を示す．E は，コア 0 の読み出しが他の
  コピーを見つけなかったので，手順 2 の無効化を省く．O は，手順 3 と 5 の
  主記憶への書き込みを，コア 0 がいずれ X を追い出すときの1回の書き込みを
  代償として省く．
\end{example}

\begin{table}[htbp]
  \centering\small
  \caption{4つのプロトコルのもとでのブロック X（\cref{ex:l19-protocols}）．
    各アクセス後のコア 0 とコア 1 の状態．b: バス操作．m: 主記憶への
    書き込み．}
  \label{tab:l19-protocols}
  \setlength{\tabcolsep}{4pt}
  \begin{tabular}{@{}cl cc cc cc cc@{}}
    \toprule
    & & \multicolumn{2}{c}{MSI} & \multicolumn{2}{c}{MOSI}
      & \multicolumn{2}{c}{MESI} & \multicolumn{2}{c}{MOESI} \\
    \cmidrule(lr){3-4}\cmidrule(lr){5-6}\cmidrule(lr){7-8}\cmidrule(l){9-10}
    手順 & アクセス & 状態 & 操作 & 状態 & 操作 & 状態 & 操作 & 状態 & 操作 \\
    \midrule
    1 & コア 0 が X を読む  & S\,I & b    & S\,I & b   & E\,I & b    & E\,I & b \\
    2 & コア 0 が X に書く & M\,I & b    & M\,I & b   & M\,I & ---  & M\,I & --- \\
    3 & コア 1 が X を読む  & S\,S & b, m & O\,S & b   & S\,S & b, m & O\,S & b \\
    4 & コア 0 が X に書く & M\,I & b    & M\,I & b   & M\,I & b    & M\,I & b \\
    5 & コア 1 が X を読む  & S\,S & b, m & O\,S & b   & S\,S & b, m & O\,S & b \\
    \midrule
    \multicolumn{2}{@{}l}{バス操作} & & 5 & & 5 & & 4 & & 4 \\
    \multicolumn{2}{@{}l}{主記憶への書き込み}  & & 2 & & 0 & & 2 & & 0 \\
    \bottomrule
  \end{tabular}
\end{table}

\begin{keypoint}
  4つのプロトコルのいずれにおいても，主記憶より新しいブロックに責任を
  持つのは M または O のキャッシュだけであり，バス操作なしに書き込める
  のは M または E のキャッシュだけである．2つの追加は独立である．E 状態は
  共有されていないデータを読んだ後の無効化を省き，O 状態は，変更された
  ブロックが他のコアに読まれたときの主記憶への書き込みを省く．
\end{keypoint}

\section{ディレクトリ方式}
\label{sec:l19-directory}

スヌーピングは，すべてのキャッシュがすべてのバス操作を観測することを
要求するので，共有バスと切り離せない．\source{Lesson 19，動画 26--30；
H\&P §5.4；\cite{agarwal1988}．}共有バスはマシンを少数のコアに制限し
（第18章），コアがネットワークを介して1つのアドレス空間を共有するより
大きなマシン（H\&P の意味での分散共有メモリ．第18章の用語に関する注を
参照）には，そもそもバスがない．すべての読み出しと無効化をネットワーク
経由で全キャッシュへ送れば，操作そのものよりはるかに大きなコストがかかる．
またバスがなければ，あるブロックに対する異なるコアの操作を順序付けるもの
が何もない．そのため，より大きなマシンは，各ブロックがどこにキャッシュ
されているかを明示的に管理する．\margin{バスが扱えるのは数個のコアだが，今日の
サーバ用ソケットは 64--128 個のコアを載せ，オンチップの
ネットワークで接続する．}

\begin{definition}[ディレクトリ方式]
  \term[でぃれくとりほうしき]{ディレクトリ方式}%
  [directory-based coherence]では，\emph{ディレクトリ}が，すべての
  ブロックについて，どのキャッシュがそのコピーを保持しているかを記録する．
  キャッシュは，スヌーピングのプロトコルならブロードキャストする操作を
  ディレクトリへ送り，ディレクトリはそれを，そのブロックを保持する
  キャッシュだけへ転送する．あるブロックに対するそうした要求はすべて
  そのディレクトリエントリを通るので，エントリが要求を直列化する．
\end{definition}

ディレクトリエントリに必要な情報はわずかである（\cref{fig:l19-directory}）．
\begin{itemize}
  \item 1つのキャッシュがそのブロックを M または E で保持しており，誰にも
    知らせずに書き込めることを示す\emph{排他}ビット．\margin{存在ビットは
    コアごと・ブロックごとに1ビットである．64 コアで 64 バイトのブロック
    なら，完全なディレクトリは主記憶の約 \SI{13}{\percent} を占める
    （第22章）．}
  \item そのキャッシュが有効なコピーを保持しているかどうかを示す，
    キャッシュごとに1つの\emph{存在}ビット．
\end{itemize}
状態そのものは依然としてキャッシュが保持するので，前節までのプロトコルを
そのまま使える．プロトコルが局所的に扱うアクセス---すべての有効な状態での
読み出しと，M および E での書き込み---はディレクトリを介さず，それ以外が
ディレクトリへ送られる．
\begin{itemize}
  \item \textbf{読み出しミス．}どのキャッシュもそのブロックを保持して
    いなければ，主記憶がブロックを供給し，読み出したコアが唯一の保持者と
    なる（MESI と MOESI では E で，排他ビットが立つ）．そうでなければ，
    ディレクトリは読み出しを，そのブロックを保持するキャッシュの1つへ
    転送し，そのキャッシュがキャッシュ間転送でブロックを供給する．
    エントリが排他であれば，読み出しは排他的な保持者へ送られ，その保持者は
    スヌープした読み出しの場合と同じように状態を変え，排他ビットが
    下ろされる．いずれの場合も，読み出したコアがエントリに加えられる．%
    \sidenote{ディレクトリはマルチコアより古い．1990 年代初頭の
    スタンフォード大学の DASH は，市販のワークステーションに
    ディレクトリボードを加えて最大 64 プロセッサのコヒーレントなマシンを
    作り，SGI Origin 2000 につながった．}
  \item \textbf{S または O での書き込み，および書き込みミス．}ディレクトリ
    は，存在ビットが立っているキャッシュだけへ無効化を転送し，書き込んだ
    コアを排他的な保持者として記録する．書き込みミスでは，ブロックを
    保持するキャッシュが I に入る前にブロックを供給し，どのキャッシュも
    保持していなければ主記憶が供給する．
\end{itemize}
キャッシュはブロックを追い出すとき---M または O では書き戻しながら---
ディレクトリに通知し，ディレクトリはそのキャッシュの存在ビットと排他
ビットを下ろす．これによりエントリは正確に保たれる．

ディレクトリは1か所に置かれる必要はない．ディレクトリは複数のスライスに
分割されて異なるノードに保持され，各スライスはブロックの集合を担当する
ので，どの1つのノードもすべての要求を処理せずに済む．あるブロックに
対する要求はすべて，そのブロックを担当するスライス，すなわちブロックの
ホームへ送られ，このスライスがそのブロックへのすべてのアクセスの順序を
決める．

\begin{figure}[htbp]
  \centering
  \resizebox{\linewidth}{!}{% Directory-based coherence: the directory entry of block X records whether
% one cache holds X exclusively and which caches hold a copy.  A write by
% core 3 is sent to the directory, which forwards the invalidation only to
% the caches that hold X.
\begin{tikzpicture}[x=1cm,y=1cm, font=\footnotesize\sffamily, >={Stealth[length=1.8mm]},
  core/.style={draw, fill=ln-accent!15, minimum width=2.1cm, minimum height=0.55cm, inner sep=2pt},
  cache/.style={draw, fill=ln-box, minimum width=2.1cm, minimum height=0.55cm, inner sep=2pt},
  cell/.style={draw, fill=white, minimum width=0.55cm, minimum height=0.5cm, inner sep=0pt},
  hd/.style={font=\scriptsize, text=ln-muted},
  lab/.style={font=\scriptsize, inner sep=2pt},
  msg/.style={->, dashed, ln-accent, thick}]
  \foreach \i/\s in {0/S, 1/O, 2/I, 3/S} {
    \node[core] (c\i) at (\i*2.8,0) {\iflangja{コア}{core}~\i};
    \node[cache] (k\i) at (\i*2.8,-0.75) {\iflangja{キャッシュ}{cache}: X (\s)};
    \draw (c\i.south) -- (k\i.north);
    \draw (k\i.south) -- (\i*2.8,-1.55);
  }
  \draw[line width=1.4pt] (-1.15,-1.55) -- (9.55,-1.55)
    node[right, font=\scriptsize, inner sep=2pt] {\iflangja{ネットワーク}{network}};
  % directory entry of X
  \node[font=\footnotesize] at (2.85,-3.1) {X};
  \foreach \j/\h/\v in {0/Ex/0, 1/C0/1, 2/C1/1, 3/C2/0, 4/C3/1} {
    \node[cell] (e\j) at (3.4+\j*0.55,-3.1) {\v};
    \node[hd, above] at (e\j.north) {\h};
  }
  \fill[ln-accent!15] (e0.south west) rectangle (e0.north east);
  \node[font=\footnotesize] at (e0) {0};
  \draw (e0.south west) rectangle (e0.north east);
  \node[draw=ln-rule, rounded corners=2pt, fit=(e0)(e4), inner xsep=16pt, inner ysep=14pt,
        yshift=2pt, xshift=-5pt] (dir) {};
  \node[lab, text=ln-muted, anchor=north] at (dir.south) {\iflangja{ディレクトリ}{directory}};
  % messages for a write by core 3
  \draw[msg, rounded corners=4pt] (k3.south) ++(0.35,0) |- node[lab, pos=0.25, right]
    {\iflangja{(1) X に書き込み}{(1) write X}} (dir.east);
  \draw[msg, rounded corners=4pt] (dir.west) -| node[lab, pos=0.75, left]
    {\iflangja{(2) 無効化}{(2) invalidate}} ([xshift=0.35cm]k0.south);
  \draw[msg] (dir.north -| k1.south) ++(0.35,0) -- node[lab, pos=0.2, right]
    {\iflangja{(2) 無効化}{(2) invalidate}} ([xshift=0.35cm]k1.south);
\end{tikzpicture}
}
  \caption{\cref{ex:l19-directory} の手順 5 におけるディレクトリ方式．
    X のエントリは，コア 0，1，3 がコピーを保持しており，どのコアも X を
    排他的には保持していないことを示す．コア 3 の書き込み (1) は，コア 0
    と 1 だけへ無効化として転送される (2)．}
  \label{fig:l19-directory}
\end{figure}

\needspace{4\baselineskip}
\begin{example}
  \label{ex:l19-directory}
  4つのコアが，ディレクトリとともに MOESI を用いる
  （\cref{tab:l19-directory}）．コア 1 が X を読み，続いて書く．その
  読み出しは他のコピーを見つけなかったので，X はコア 1 のキャッシュで E に
  あり，書き込みはメッセージを必要としない．コア 3 の読み出しは
  ディレクトリへ行き，ディレクトリはそれを排他的な保持者であるコア 1 へ
  転送する．コア 1 はブロックを供給し，その所有者となる．ディレクトリは
  コア 0 の読み出しを保持者の1つ，やはりコア 1 へ転送する．手順 5 で
  コア 3 が X に書くと，ディレクトリはコア 0 と 1 だけへ無効化を送る．
  コピーを保持していないコア 2 にはメッセージが届かない．ブロード
  キャストであればコア 2 にも届いていたはずである．
\end{example}

\begin{table}[htbp]
  \centering\small
  \caption{ディレクトリ方式の MOESI のもとでのブロック X
    （\cref{ex:l19-directory}）．各アクセス後のキャッシュの状態と
    ディレクトリエントリ（排他ビット，コア 0 から 3 の存在ビット），
    および送られたメッセージ（d: ディレクトリ）．}
  \label{tab:l19-directory}
  \setlength{\tabcolsep}{4pt}
  \begin{tabular}{@{}cl cccc cc l@{}}
    \toprule
    & & \multicolumn{4}{c}{コアのキャッシュ} & \multicolumn{2}{c}{エントリ} & \\
    \cmidrule(lr){3-6}\cmidrule(lr){7-8}
    手順 & アクセス & 0 & 1 & 2 & 3 & 排他 & 存在 & メッセージ \\
    \midrule
    1 & コア 1 が X を読む  & I & E & I & I & 1 & 0100 & 1$\to$d，主記憶$\to$1 \\
    2 & コア 1 が X に書く & I & M & I & I & 1 & 0100 & --- \\
    3 & コア 3 が X を読む  & I & O & I & S & 0 & 0101 & 3$\to$d$\to$1$\to$3 \\
    4 & コア 0 が X を読む  & S & O & I & S & 0 & 1101 & 0$\to$d$\to$1$\to$0 \\
    5 & コア 3 が X に書く & I & I & I & M & 1 & 0001 & 3$\to$d$\to$0，1 \\
    6 & コア 3 が X に書く & I & I & I & M & 1 & 0001 & --- \\
    \bottomrule
  \end{tabular}
\end{table}

\section{キャッシュミスとコヒーレンス}
\label{sec:l19-misses}

第14章では，キャッシュミスを初期参照ミス，容量性ミス，競合性ミスに分けた．
無効化は第4の種類を加える．\source{Lesson 19，動画 31--33；H\&P §5.3．}

\begin{definition}[コヒーレンスミス]
  \term[こひーれんすみす]{コヒーレンスミス}[coherence miss]とは，他の
  コアがブロックに書いたためにキャッシュ内のコピーが無効化された，その
  ブロックに対するミスである．
\end{definition}

コヒーレンスミスには2つの種類がある．
\begin{description}
  \item[真共有] コアは，他のコアが実際に書いたデータにアクセスする．
    このような\term[しんきょうゆう]{真共有}[true sharing]では，ミスは
    必要である．ミスがなければ，コアは古い値を読んでしまうからである．
  \item[偽共有] 他のコアは，たまたま同じブロックにある別のデータに
    書いた．コアがアクセスするデータは変化しておらず，ミスは，
    コヒーレンスがアドレス単位ではなくブロック単位で保たれていることだけ
    に起因する．このような\term[ぎきょうゆう]{偽共有}[false sharing]は
    ブロックサイズとともに増える．大きなブロックほど無関係なデータを
    多く含むからである．
\end{description}

\begin{example}
  \label{ex:l19-false-sharing}
  コア 0 とコア 1 のスレッドが，それぞれ自分のカウンタ c0 と c1 を
  \num{1000} 回インクリメントし，2つのスレッドのインクリメントが交互に
  起こるとする．2つのカウンタは隣接しており，同じキャッシュブロックに
  ある（\cref{fig:l19-false-sharing}）．MOESI のもとでは，コア 0 の最初の
  インクリメントは他のコピーを見つけないので，その書き込みは無効化を
  必要としない．しかしコア 1 の最初のインクリメント以降は，どの
  インクリメントも，他方のコアがコピーを保持している状態でブロックに書き，
  それを無効化する．そのため，コア 0 の2回目以降のインクリメントはすべて
  コヒーレンスミスから始まる．どちらのスレッドも相手のカウンタを読む
  ことはないので，\num{1998} 回のミスはすべて偽共有によるものである．
  カウンタが別々のブロックにあれば，残るのは2回の初期参照ミスだけで
  ある．\tip{異なるスレッドが書くデータは切り離す．x86-64 では
  \texttt{alignas(64)} を使う．Linux の \texttt{perf c2c} は，それが
  できていないラインを見つける．}
\end{example}
\begin{marginfigure}
  \resizebox{\linewidth}{!}{% False sharing: two counters used by different cores lie in one cache block.
\begin{tikzpicture}[x=1cm,y=1cm, font=\footnotesize\sffamily, >={Stealth[length=1.8mm]},
  core/.style={draw, fill=ln-accent!15, minimum width=1.5cm, minimum height=0.55cm, inner sep=2pt},
  cell/.style={draw, minimum width=0.5cm, minimum height=0.5cm, inner sep=0pt}]
  \fill[ln-box] (-0.25,-0.25) rectangle (0.75,0.25);
  \foreach \j in {0,...,7} {
    \node[cell] (b\j) at (\j*0.5,0) {};
  }
  \node[font=\scriptsize] at (b0) {c0};
  \node[font=\scriptsize] at (b1) {c1};
  \node[core] (k0) at (0,1.2)   {\iflangja{コア}{core}~0};
  \node[core] (k1) at (0.5,-1.2) {\iflangja{コア}{core}~1};
  \draw[->] (k0.south) -- (b0.north);
  \draw[->] (k1.north) -- (b1.south);
  \node[font=\scriptsize, text=ln-muted, anchor=north east] at (b7.south east)
    {\iflangja{1つのキャッシュブロック}{one cache block}};
\end{tikzpicture}
}
  \caption{偽共有: 異なるコアが使うカウンタが同じキャッシュブロックに
    ある．}
  \label{fig:l19-false-sharing}
\end{marginfigure}

\begin{keypoint}[章のまとめ]
  共有メモリ型マルチプロセッサの専用キャッシュはコヒーレントでなければ
  ならない．すなわち，コアは自分の書き込みを見，他のコアの書き込みを短い
  時間の後に見，ある位置へのすべての書き込みを他のどのコアとも同じ順序で
  見る．共有バスでは，キャッシュどうしが互いの操作をスヌープし，バスが
  それらを順序付ける．書き込み時更新方式はすべての書き込みをすべての
  コピーへブロードキャストし，書き込み時無効化方式はブロックごとに1回の
  無効化をブロードキャストして，一般的なアクセスパターンでは後者が勝る．
  MSI は3つの状態で書き込み時無効化方式を実現する．キャッシュ間転送は
  現在の値を保持するキャッシュから読み出しに応答し，MOSI の O 状態は転送
  のたびにその値を主記憶へ書くことを避け，MESI の E 状態は1つのコアだけが
  使うデータに対する無効化を避ける．MOESI は両方を組み合わせる．共有バス
  を持たないマシンは，ブロードキャストを，どのキャッシュが各ブロックを
  保持しているかを記録し，そのブロックへの要求を順序付けるディレクトリで
  置き換える．無効化は，ミスの第4の種類であるコヒーレンスミスを引き起こし，
  これは真共有または偽共有に由来する．
\end{keypoint}

\end{document}
